\section{Analiza SWOT}

\subsection{Puncte tari (\textit{Strengths})}
\begin{itemize}
    \item \textbf{High-value use case}: industria farmaceutică gestionează o nevoie umană critică și recurentă. Întrebări despre medicamente sunt puse zilnic, iar existența unei soluții care să eficientizeze timpul de interacțiune al farmacistului cu clienții este esențială.
    \item \textbf{Agentic Workflow}: utilizarea mai multor agenți specializați în favoarea unui singur agent generic crește acuratețea răspunsurilor și scade riscul de halucinare.
    \item \textbf{Flexible monetization}: farmaciile independente sau mai puțin dezvoltate pot integra soluția prin planul \textit{Starter} accesibil (de la 10\,\$/lună) și prin pachete suplimentare de mesaje \textit{pay-as-you-go}, fără un angajament financiar mare la început.
    \item \textbf{Dual audience capability}: farmaciile pot pune serviciul la dispoziția clienților prin contul de farmacist; clienții se pot informa și individual despre medicamentele pe care le consumă.
    \item \textbf{Medication normalization engine \& cross-country drug alternatives}: sistemul identifică medicamente și ingrediente în pofida greșelilor de scriere, oferă alternative locale și echivalente disponibile în alte țări.
\end{itemize}

\subsection{Puncte slabe (\textit{Weaknesses})}
\begin{itemize}
    \item \textbf{Data dependency}: agentul de \textit{facts} are nevoie de baze de date cu medicamente și ingrediente care să conțină informații corecte și să fie actualizate periodic.
    \item \textbf{Bariera de adoptare}: farmaciștii sunt presați de timp și lucrează cu volume mari de clienți zilnic; pot fi reticenți în a integra un instrument suplimentar dacă acesta nu este integrat în sistemul de gestiune existent.
    \item \textbf{High liability risk}: halucinațiile legate de dozaje sau interacțiuni medicamentoase, însoțite de neatenție umană, pot pune în pericol viața consumatorilor și pot atrage consecințe legale.
\end{itemize}

\subsection{Oportunități (\textit{Opportunities})}
\begin{itemize}
    \item \textbf{Expansiune către industria cosmetică}: soluția poate fi extinsă pentru a explica ingredientele din produsele cosmetice și potențialele riscuri adverse.
    \item \textbf{Integrare ERP / Gestiune}: soluția poate fi integrată cu softurile de gestiune farmaceutică și poate oferi alternative de medicamente bazate pe stocurile existente.
    \item \textbf{Internaționalizare}: motorul de echivalențe trans-frontaliere deschide piețe în care turismul medical sau migrația cresc cererea pentru identificarea de alternative locale.
\end{itemize}

\subsection{Amenințări (\textit{Threats})}
\begin{itemize}
    \item \textbf{Reglementări stricte din domeniul medical}: pot restricționa sau elimina utilizarea AI în consultanță medicală (ex.: AI Act în UE).
    \item \textbf{Negative perception of AI in medical contexts}: atât farmaciștii, cât și clienții, pot fi sceptici în a adopta o soluție bazată pe AI când vine vorba de sănătate, preferând să se bazeze pe specialiști.
    \item \textbf{Concurența Big Pharma}: companiile mari ar putea intra pe piață cu propriile soluții și ar putea fi preferate datorită renumelui deja stabilit.
\end{itemize}
